%! root=./main.tex
\section{Related Works}

CCS'20 survey on efficacy of CFI:
https://dl-acm-org.libproxy.utdallas.edu/doi/pdf/10.1145/3372297.3417867
* Real boundaries differ from claimed boundaries due to implementation or algorithmic errors
* Five of seven common mistakes were in the implementation stage
    * unexpected optimization
    * corner-cases: inlined assembly, dynamic shared objects (DSO)
    * logic bugs in implementation
    * mismatched specification, e.g., QEMU bug in ARM PA emulation, not setting the NX bit
    * adding unintended permissible jump targets (xend bytes in TSX-based CFI)
* Mentioned binary-compatible CFI:  Opaque CFI [ 37 ], CCFIR [ 72 ], binCFI [ 75] and CFIMon 
* Mentioned theoretically-only safe CFI:  SafeDispatch [ 28 ] and PARTS [ 34 ]
* CScan tool very limited: relies on source-level analysis, only tests each ICT guard code once
* Mention false positives, maybe we can mention we can analyze the implementation and prove it precludes false positives
* Analyzed 56 CFI solution evaluation methods, found three methods:
    1) theoretical evaluation
    2) experimental verification
    3) quantitative analysis

Franz'17 survey on CFI:
https://dl.acm.org/doi/pdf/10.1145/3054924
* common security metric AIR is approximate and leaves holes for attackers
* aside: Carlini et al Control Flow Bending show even with the most precise CFG and a shadowstack attacks may still be possible

ARM BTI binary rewriter:
10.1109/TIFS.2025.3580342
* does not even try to make binary safe, just make it harder to attack
* binary analysis, incremental rewriting, runtime monitor for distinguishing attacks from ?

